\documentclass[12pt,a4paper]{article}
\usepackage[UTF8]{ctex}
\usepackage{amsmath,amssymb,amsthm}
\usepackage{geometry}

\geometry{margin=2.5cm}

\title{向量叉乘与矩阵乘法的变换形式：$a \times b \cdot c$ 的等价表达式}
\author{第8章补充说明}
\date{}

\begin{document}

\maketitle

\section{问题提出}

\textbf{问题}：表达式 $a \times b \cdot c$（或 $[a]_{\times} b \cdot c$）有哪些等价的变换形式？

\textbf{符号说明}：
\begin{itemize}
\item $a, b, c \in \mathbb{R}^3$：三维向量
\item $a \times b$：向量叉乘
\item $[a]_{\times}$：向量$a$的反对称矩阵
\item $\cdot$ 或矩阵乘法：向量点积或矩阵乘法
\end{itemize}

\section{基本形式}

\subsection{形式1：直接叉乘形式}

\begin{equation}
(a \times b) \cdot c
\end{equation}

这是最直接的形式，先计算叉乘，再与$c$做点积。

\subsection{形式2：反对称矩阵形式}

使用反对称矩阵表示叉乘：
\begin{equation}
[a]_{\times} b \cdot c = (a \times b) \cdot c
\end{equation}

其中反对称矩阵定义为：
\begin{equation}
[a]_{\times} = \begin{bmatrix}
0 & -a_z & a_y \\
a_z & 0 & -a_x \\
-a_y & a_x & 0
\end{bmatrix}
\end{equation}

\textbf{性质}：
\begin{equation}
a \times b = [a]_{\times} b = -[b]_{\times} a
\end{equation}

\section{主要变换形式}

\subsection{变换1：标量三重积形式}

\textbf{标量三重积}：
\begin{equation}
(a \times b) \cdot c = a \cdot (b \times c) = b \cdot (c \times a)
\end{equation}

\textbf{证明}：
\begin{align}
(a \times b) \cdot c &= \det([a, b, c]) \\
&= a \cdot (b \times c) \\
&= b \cdot (c \times a) \\
&= c \cdot (a \times b)
\end{align}

\textbf{等价形式}：
\begin{align}
(a \times b) \cdot c &= a \cdot (b \times c) \\
&= b \cdot (c \times a) \\
&= c \cdot (a \times b) \\
&= -(b \times a) \cdot c \\
&= -(a \times c) \cdot b \\
&= -(c \times b) \cdot a
\end{align}

\subsection{变换2：反对称矩阵的转置}

\textbf{反对称矩阵的性质}：
\begin{equation}
[a]_{\times}^T = -[a]_{\times}
\end{equation}

\textbf{应用}：
\begin{align}
[a]_{\times} b \cdot c &= b^T [a]_{\times}^T c \\
&= -b^T [a]_{\times} c \\
&= -b \cdot (a \times c)
\end{align}

\textbf{另一种形式}：
\begin{align}
[a]_{\times} b \cdot c &= ([a]_{\times} b)^T c \\
&= b^T [a]_{\times}^T c \\
&= -b^T [a]_{\times} c \\
&= -b \cdot (a \times c)
\end{align}

\subsection{变换3：矩阵乘法的结合律}

\textbf{形式A}：
\begin{equation}
[a]_{\times} b \cdot c = ([a]_{\times} b)^T c = b^T [a]_{\times}^T c
\end{equation}

\textbf{形式B}：
\begin{equation}
[a]_{\times} b \cdot c = b^T ([a]_{\times}^T c) = b^T (-[a]_{\times} c) = -b \cdot (a \times c)
\end{equation}

\textbf{形式C}（使用标量三重积）：
\begin{equation}
[a]_{\times} b \cdot c = (a \times b) \cdot c = a \cdot (b \times c) = a^T [b]_{\times} c
\end{equation}

\subsection{变换4：旋转矩阵的伴随性质}

\textbf{关键关系}：对于旋转矩阵$R \in SO(3)$和向量$p$：
\begin{equation}
R [p]_{\times} = [Rp]_{\times} R
\end{equation}

\textbf{证明}：
对于任意向量$q$：
\begin{align}
R [p]_{\times} q &= R (p \times q) \\
&= (Rp) \times (Rq) \quad \text{（旋转保持叉积）} \\
&= [Rp]_{\times} (Rq) \\
&= [Rp]_{\times} R q
\end{align}

\textbf{应用}：
\begin{align}
R (a \times b) \cdot c &= (Ra) \times (Rb) \cdot c \\
&= [Ra]_{\times} (Rb) \cdot c \\
&= ([Ra]_{\times} R b) \cdot c \\
&= (R [a]_{\times} b) \cdot c
\end{align}

\textbf{等价形式}：
\begin{align}
R [a]_{\times} b \cdot c &= [Ra]_{\times} R b \cdot c \\
&= (Ra) \times (Rb) \cdot c \\
&= R (a \times b) \cdot c
\end{align}

\subsection{变换5：转置形式}

\textbf{对于$R^T$}：
\begin{equation}
R^T [p]_{\times} = [R^T p]_{\times} R^T
\end{equation}

\textbf{应用}：
\begin{align}
R^T (a \times b) \cdot c &= (R^T a) \times (R^T b) \cdot c \\
&= [R^T a]_{\times} (R^T b) \cdot c \\
&= ([R^T a]_{\times} R^T b) \cdot c \\
&= (R^T [a]_{\times} b) \cdot c
\end{align}

\section{完整的变换形式总结}

\subsection{基本等价形式（共6种）}

\begin{enumerate}
\item \textbf{直接形式}：
\begin{equation}
(a \times b) \cdot c
\end{equation}

\item \textbf{反对称矩阵形式}：
\begin{equation}
[a]_{\times} b \cdot c
\end{equation}

\item \textbf{标量三重积形式1}：
\begin{equation}
a \cdot (b \times c)
\end{equation}

\item \textbf{标量三重积形式2}：
\begin{equation}
b \cdot (c \times a)
\end{equation}

\item \textbf{转置形式}：
\begin{equation}
b^T [a]_{\times}^T c = -b^T [a]_{\times} c = -b \cdot (a \times c)
\end{equation}

\item \textbf{矩阵形式}：
\begin{equation}
a^T [b]_{\times} c
\end{equation}
\end{enumerate}

\subsection{带旋转矩阵的变换形式}

\begin{enumerate}
\item \textbf{旋转后叉乘}：
\begin{equation}
(Ra) \times (Rb) \cdot c
\end{equation}

\item \textbf{旋转矩阵与反对称矩阵}：
\begin{equation}
R [a]_{\times} b \cdot c = [Ra]_{\times} R b \cdot c
\end{equation}

\item \textbf{转置旋转矩阵}：
\begin{equation}
R^T [a]_{\times} b \cdot c = [R^T a]_{\times} R^T b \cdot c
\end{equation}
\end{enumerate}

\section{具体应用示例}

\subsection{示例1：伴随变换中的应用}

在伴随变换中，我们经常遇到：
\begin{equation}
R [p]_{\times} \omega \cdot v
\end{equation}

\textbf{变换形式}：
\begin{align}
R [p]_{\times} \omega \cdot v &= [Rp]_{\times} R \omega \cdot v \\
&= (Rp) \times (R \omega) \cdot v \\
&= R (p \times \omega) \cdot v
\end{align}

\subsection{示例2：速度合成中的应用}

在速度合成中：
\begin{equation}
R^T [p]_{\times} \omega \cdot v
\end{equation}

\textbf{变换形式}：
\begin{align}
R^T [p]_{\times} \omega \cdot v &= [R^T p]_{\times} R^T \omega \cdot v \\
&= (R^T p) \times (R^T \omega) \cdot v \\
&= R^T (p \times \omega) \cdot v
\end{align}

\section{重要恒等式}

\subsection{恒等式1：反对称矩阵的转置}

\begin{equation}
[a]_{\times}^T = -[a]_{\times}
\end{equation}

\subsection{恒等式2：旋转保持叉积}

\begin{equation}
R(a \times b) = (Ra) \times (Rb)
\end{equation}

\subsection{恒等式3：反对称矩阵与旋转矩阵}

\begin{equation}
R [a]_{\times} = [Ra]_{\times} R
\end{equation}

\begin{equation}
R^T [a]_{\times} = [R^T a]_{\times} R^T
\end{equation}

\subsection{恒等式4：标量三重积的循环性}

\begin{equation}
(a \times b) \cdot c = (b \times c) \cdot a = (c \times a) \cdot b
\end{equation}

\subsection{恒等式5：反对称矩阵的反对称性}

\begin{equation}
[a]_{\times} b = -[b]_{\times} a
\end{equation}

\section{总结}

\textbf{主要变换形式}：

对于表达式 $(a \times b) \cdot c$ 或 $[a]_{\times} b \cdot c$，有以下主要变换形式：

\begin{enumerate}
\item \textbf{标量三重积形式}：$a \cdot (b \times c) = b \cdot (c \times a) = c \cdot (a \times b)$
\item \textbf{反对称矩阵转置形式}：$b^T [a]_{\times}^T c = -b^T [a]_{\times} c = -b \cdot (a \times c)$
\item \textbf{矩阵形式}：$a^T [b]_{\times} c$
\item \textbf{旋转矩阵形式}：$R [a]_{\times} b \cdot c = [Ra]_{\times} R b \cdot c$
\item \textbf{转置旋转矩阵形式}：$R^T [a]_{\times} b \cdot c = [R^T a]_{\times} R^T b \cdot c$
\end{enumerate}

\textbf{应用场景}：
\begin{itemize}
\item 机器人学中的速度合成
\item 伴随变换的推导
\item 动力学方程的简化
\item 约束力的计算
\end{itemize}

\end{document}

